% Section 6: Pointwise VPC Inference
% Requires amsmath, amssymb, amsthm.
% Uses the notation and theorem environments of Sections 2--5.

\section{Pointwise VPC Significance Testing}
\label{sec:pointwise_inference}


The omnibus hypothesis $\Gamma=0$ rules out
random-effects variation at all covariate values.
In contrast, the pointwise null permits variation
in random-effects directions that do not contribute
at a specified covariate value. This distinction
changes the null geometry from a vertex to a ray
and requires likelihood-ratio calibration that
accounts for both the vertex and nonzero points
of the null set.

The general theory of likelihood inference under
boundary constraints provides the local
tangent-cone framework
\citep{self1987asymptotic,shapiro1985}.
We first establish fixed-parameter limits and
then examine local null sequences approaching
the vertex. An endpoint-envelope calibration is established
through the corresponding Gaussian projection
experiment.

\subsection{Pointwise hypotheses and their covariance geometry}
\label{subsec:pointwise_null}

Fix $x_0\in\mathbb R$, and put
\[
z_0=(1,x_0)^{\mathsf T},\qquad
w_0=(-x_0,1)^{\mathsf T}.
\]
By Proposition~\ref{prop:zero_vpc}, provided the
observation-level variance is positive,
\begin{equation}\label{eq:sec6-pointwise-null}
\begin{aligned}
H_0(x_0):\rho(x_0;\theta)=0
\quad\Longleftrightarrow\quad
\Gamma z_0=0.
\end{aligned}
\end{equation}
For a positive-semidefinite $2\times2$ matrix this is
equivalent to
\begin{equation}\label{eq:sec6-null-ray}
\begin{aligned}
K_0(x_0)=\bigl\{\gamma:\ &
\Gamma(\gamma)=t w_0w_0^{\mathsf T},\\
& t\geq0\bigr\}.
\end{aligned}
\end{equation}
Write $r_0=\operatorname{vech}(w_0w_0^{\mathsf T})$
using the covariance-coordinate convention of
Section~\ref{sec:geometry_asymptotics},
so that $K_0(x_0)=\mathbb R_+r_0$. The full null parameter
space is $\Theta_0(x_0)=\Lambda\times K_0(x_0)$ and the
alternative parameter space is $\Theta=\Lambda\times K$.
The pointwise null is therefore composite: it contains
the covariance vertex and nonzero rank-one matrices.
For a binary covariate, the separate nulls at $x_0=0$
and $x_0=1$ are distinct; their conjunction is equivalent
to $\Gamma=0$.

\subsection{Pointwise likelihood-ratio statistic}
\label{subsec:pointwise_lrt}

Let $\widehat\theta$ maximize the marginal likelihood over
$\Theta$, and let $\widetilde\theta_{x_0}$ maximize it over
$\Theta_0(x_0)$. Define
\begin{equation}\label{eq:sec6-lrt}
\begin{aligned}
\mathrm{LRT}_n(x_0)
&=2\bigl\{\ell_n(\widehat\theta)\\
&\hspace{1.6em}
-\ell_n(\widetilde\theta_{x_0})\bigr\}.
\end{aligned}
\end{equation}
The null fit uses one nonnegative variance parameter $t$
in $\Gamma=t w_0w_0^{\mathsf T}$, while all fixed effects
and family-specific parameters remain unrestricted within
$\Lambda$. In particular, the null is not imposed by setting
$\Gamma=0$ or by applying a Wald approximation to
$\widehat\rho(x_0)$.

\subsection{Tangent geometry and nested-cone limit}
\label{subsec:pointwise_tangent}

For $\gamma_0=t_0r_0$, the null tangent cone is
\begin{equation}\label{eq:sec6-null-tangent}
T_{K_0(x_0)}(\gamma_0)=
\begin{cases}
\operatorname{span}(r_0),&t_0>0,\\
K_0(x_0),&t_0=0.
\end{cases}
\end{equation}
The alternative tangent cone is the rank-one supporting
half-space of
Proposition~\ref{prop:sec4-tangent-cones}
when $t_0>0$, and $K$ when $t_0=0$.
Both full tangent cones have the common unrestricted
nuisance factor $\mathbb R^{p_\lambda}$.

Let $I_0=I(\theta_0)$ denote the normalized full Fisher
information and $J_0$ its efficient covariance
information, defined by the Schur complement in
\eqref{eq:sec4-efficient-information}.

\begin{theorem}\label{thm:sec6-nested-cones}
Suppose $H_0(x_0)$ holds at a fixed $\theta_0$,
$J_0\succ0$, and the local asymptotic normality,
local-set convergence, consistency, and constrained-argmax
conditions of Section~\ref{sec:geometry_asymptotics}
hold for both constrained fits.
Put $T_1=T_K(\gamma_0)$ and
$T_0=T_{K_0(x_0)}(\gamma_0)$, and set
$C_j=J_0^{1/2}T_j$. Then
\begin{equation}\label{eq:sec6-projection-limit}
\begin{aligned}
\mathrm{LRT}_n(x_0)
&\xrightarrow{d}
\|\Pi_{C_1}(G)\|^2\\
&\quad-\|\Pi_{C_0}(G)\|^2,\\
G&\sim N_3(0,I_3).
\end{aligned}
\end{equation}
\end{theorem}

The common nuisance coordinates disappear after profiling
the local quadratic likelihood. Unlike the omnibus vertex
test, the second projection in
\eqref{eq:sec6-projection-limit} is generally nonzero.
The proof is given in Appendix~\ref{app:pointwise}.

\subsection{Fixed nonzero rank-one and vertex nulls}
\label{subsec:pointwise_fixed_limits}

\begin{proposition}\label{prop:sec6-rankone-lrt}
Under the conditions of
Theorem~\ref{thm:sec6-nested-cones}, if
$\Gamma_0=t_0w_0w_0^{\mathsf T}$ with $t_0>0$, then
\begin{equation}\label{eq:sec6-rankone-mixture}
\mathrm{LRT}_n(x_0)\xrightarrow{d}
\tfrac12\chi_1^2+\tfrac12\chi_2^2.
\end{equation}
\end{proposition}

At a nonzero rank-one point, the null tangent is a line
contained in the boundary hyperplane of the alternative
tangent half-space. The line is orthogonal to the half-space
normal even after efficient-information whitening, so
removing the common line leaves one unrestricted Gaussian
coordinate and one one-sided coordinate. Appendix~\ref{app:pointwise-rankone} gives
the projection calculation. The mixture weights are
independent of $J_0$.

At $\Gamma_0=0$, the null tangent is instead the ray
$K_0(x_0)$ and the alternative tangent is $K$. Thus
\begin{equation}\label{eq:sec6-vertex-limit}
\begin{aligned}
T_{\mathrm v}(J_0,x_0)
&=\|\Pi_{J_0^{1/2}K}(G)\|^2\\
&\quad-
\|\Pi_{J_0^{1/2}K_0(x_0)}(G)\|^2.
\end{aligned}
\end{equation}

The vertex distribution admits an exact spherical
integral representation. In particular, its CDF,
density on $(0,\infty)$, and possible atom at zero
are given in Appendix~\ref{app:pointwise-vertex-density}.
These expressions provide a deterministic route
to numerical evaluation of the vertex reference
distribution without assuming a finite
chi-bar-square mixture.

This is not the omnibus statistic of
Section~\ref{sec:global_lrt} and is
not, in general, distributed according to
\eqref{eq:sec6-rankone-mixture}. We make no universal
chi-bar-square assertion for this difference of projections.

\subsection{Local null sequences approaching the vertex}
\label{subsec:pointwise_local_sequences}

The fixed-parameter limits do not, by themselves, control
null sequences whose nonzero covariance component shrinks
at the $n^{-1/2}$ scale. To describe this regime, let
$\theta_*=(\lambda_*,0)$ be a vertex parameter and let
$J_*=J(\theta_*)\succ0$. Define
\begin{equation}\label{eq:sec6-whitened-cones}
\begin{gathered}
C=J_*^{1/2}K,\qquad R=\mathbb R_+e,\\
e=\frac{J_*^{1/2}r_0}
{\|J_*^{1/2}r_0\|}.
\end{gathered}
\end{equation}
Consider null sequences
\begin{equation}\label{eq:sec6-local-sequence}
\begin{aligned}
\gamma_n
&=\frac{a\,r_0}
{\sqrt n\,\|J_*^{1/2}r_0\|},
\qquad a\geq0,
\end{aligned}
\end{equation}
with nuisance parameters converging to $\lambda_*$ under
the local experiment assumptions. Thus $a$ is the
Fisher-standardized null-ray displacement, not the
unstandardized variance $t$.

\begin{proposition}\label{prop:sec6-local-limit}
Suppose the likelihoods and constrained local parameter
sets admit a joint Gaussian-shift limit along
\eqref{eq:sec6-local-sequence}, with efficient covariance
information $J_*$. Then
\begin{equation}\label{eq:sec6-local-limit}
\begin{aligned}
\mathrm{LRT}_n(x_0)&\xrightarrow{d}T_a,\\
T_a&=\|\Pi_C(G+ae)\|^2\\
&\quad-\|\Pi_R(G+ae)\|^2,\\
G&\sim N_3(0,I_3).
\end{aligned}
\end{equation}
The endpoints satisfy
$T_0\stackrel d=T_{\mathrm v}(J_*,x_0)$
and
\begin{equation}\label{eq:sec6-infinity-limit}
\begin{gathered}
T_a\xrightarrow{d}T_\infty
\quad(a\to\infty),\\
T_\infty\sim
\tfrac12\chi_1^2+\tfrac12\chi_2^2.
\end{gathered}
\end{equation}
\end{proposition}

The Gaussian-shift representation and tangent
half-space limit are established in
Appendices~\ref{app:pointwise-local} and
\ref{app:pointwise-rankone}.
The assertion for finite $a$ concerns a local sequence
and is distinct from the fixed nonzero-rank-one result.

\subsection{Endpoint maximality of the Gaussian null tail}
\label{subsec:pointwise_endpoint}

The following geometric result supplies a single calibration
across the entire Gaussian local-null ray. It is stated in
a form that does not require the special Lorentz
representation.

\begin{theorem}[Gaussian endpoint maximality]
\label{thm:sec6-endpoint}
Let $C\subseteq\mathbb R^d$ be a closed convex cone and
let $R=\mathbb R_+e\subseteq C$, where $\|e\|=1$.
For $G\sim N_d(0,I_d)$, set
\[
\begin{aligned}
T_a={}&\|\Pi_C(G+ae)\|^2\\
&-\|\Pi_R(G+ae)\|^2,
\qquad a\geq0.
\end{aligned}
\]
For every $u>0$,
\begin{equation}\label{eq:sec6-endpoint-max}
\begin{gathered}
\sup_{a\geq0}\Pr(T_a>u)\\
=\max\Bigl\{
\Pr(T_0>u),\\
\lim_{a\to\infty}\Pr(T_a>u)
\Bigr\}.
\end{gathered}
\end{equation}
\end{theorem}

The proof in Appendix~\ref{app:pointwise-endpoint} conditions on the Gaussian
coordinates orthogonal to $e$. On the negative half-line,
the conditional acceptance profile is log-concave by
convexity and Pr\'ekopa's theorem; on the positive half-line,
it is nonincreasing by translation along $e$ inside the
cone. The resulting profile is unimodal, and Gaussian
convolution preserves this unimodality.

For the present three-dimensional whitened PSD cone,
the boundary is smooth away from the vertex, and
\eqref{eq:sec6-infinity-limit} identifies the positive-ray
endpoint. Consequently,
\begin{equation}\label{eq:sec6-pointwise-tail-bound}
\begin{gathered}
\Pr(T_a>u)\leq
\max\{p_{\mathrm v}(u),p_\infty(u)\},\\
p_{\mathrm v}(u)
=\Pr\{T_{\mathrm v}(J_*,x_0)>u\},\\
\begin{aligned}
p_\infty(u)
={}&\tfrac12\Pr(\chi_1^2>u)\\
&+\tfrac12\Pr(\chi_2^2>u).
\end{aligned}
\end{gathered}
\end{equation}
where $a\geq0$.
The endpoint bound concerns the limiting Gaussian
experiment; its transfer to a uniform statement for the
original GLMM requires additional statistical conditions.

\subsection{Calibration and feasible pointwise testing}
\label{subsec:pointwise_calibration}

Let $c_{\mathrm v,\alpha}(J,x_0)$ be an upper $\alpha$
critical value of the vertex statistic in
\eqref{eq:sec6-vertex-limit}, and let
$c_{\infty,\alpha}$ be an upper $\alpha$ critical value
of $\tfrac12\chi_1^2+\tfrac12\chi_2^2$. Define
\begin{equation}\label{eq:sec6-critical-value}
\begin{aligned}
c_{\alpha}^{\mathrm{pw}}(J,x_0)
=\max\bigl\{
&c_{\mathrm v,\alpha}(J,x_0),\\
&c_{\infty,\alpha}
\bigr\}.
\end{aligned}
\end{equation}
By Theorem~\ref{thm:sec6-endpoint}, the Gaussian experiment
satisfies
\begin{equation}\label{eq:sec6-gaussian-size}
\begin{gathered}
\sup_{a\geq0}
\Pr\bigl\{T_a>
c_{\alpha}^{\mathrm{pw}}(J_*,x_0)\bigr\}\\
\leq\alpha.
\end{gathered}
\end{equation}
For feasible endpoint-envelope calibration, let
$\widetilde\lambda_n$ be the nuisance component of the
pointwise restricted estimator and set
\begin{equation}\label{eq:sec6-vertex-information}
\widehat J_{V,n}=J(\widetilde\lambda_n,0).
\end{equation}
This vertex-evaluated information is used in the vertex
reference distribution; the remote-ray reference is
independent of $J$. Reject $H_0(x_0)$ when
\begin{equation}\label{eq:sec6-feasible-test}
\mathrm{LRT}_n(x_0)>
c_{\alpha}^{\mathrm{pw}}(\widehat J_{V,n},x_0).
\end{equation}
The vertex quantile can be evaluated by Gaussian projection
simulation or deterministic integration over the transformed
cone. Its calculation must use the difference of two squared
projection norms, not the omnibus chi-bar-square weights
of Section~\ref{sec:global_lrt}.

At the vertex and along local null sequences, consistency
of $\widetilde\lambda_n$ for the vertex nuisance parameter
makes \eqref{eq:sec6-vertex-information} appropriate. At a
fixed nonzero null covariance, the limiting LRT is the
remote-ray mixture, whose cutoff does not depend on the
information matrix. Uniform validity across both regimes,
however, requires a stronger approximation than pointwise
likelihood asymptotics alone.

\begin{theorem}[Conditional uniform calibration]
\label{thm:sec6-uniform-size}
Fix $x_0\in\mathbb R$, $0<\alpha<1$, and a compact
null class
\[
\Theta_0^T=\{(\lambda,t r_0):
\lambda\in\Lambda_0,\ 0\le t\le T\},\qquad T<\infty.
\]
Assume $J_V(\lambda):=J(\lambda,0)$ is continuous and
positive definite on a neighborhood of $\Lambda_0$.
For $\theta=(\lambda,t r_0)\in\Theta_0^T$, put
\begin{equation}\label{eq:sec6-uniform-displacement}
\begin{aligned}
 a_n(\theta)&=\sqrt n\,t\,
 \|J_V(\lambda)^{1/2}r_0\|,\\
 e_\lambda&=\frac{J_V(\lambda)^{1/2}r_0}
 {\|J_V(\lambda)^{1/2}r_0\|}.
\end{aligned}
\end{equation}
Let $T_a^{J_V(\lambda)}$ denote the Gaussian
cone-versus-ray statistic in
\eqref{eq:sec6-local-limit}, with $J_*$ replaced by
$J_V(\lambda)$ and displacement $a$; at $a=\infty$
use the remote-ray limit in
\eqref{eq:sec6-infinity-limit}.

Suppose the restricted nuisance estimator and the
vertex-information plug-in are well defined and the
following \emph{uniform rejection-probability approximation}
holds for the feasible test in
\eqref{eq:sec6-feasible-test}:
\begin{equation}\label{eq:sec6-uniform-approx}
\begin{aligned}
\sup_{\theta\in\Theta_0^T}\Bigl|
&P_\theta\{\mathrm{LRT}_n(x_0)>
 c_\alpha^{\mathrm{pw}}(\widehat J_{V,n},x_0)\}\\
-&P\{T_{a_n(\theta)}^{J_V(\lambda)}>
 c_\alpha^{\mathrm{pw}}(J_V(\lambda),x_0)\}
\Bigr|\longrightarrow0.
\end{aligned}
\end{equation}
Since $c_{\infty,\alpha}>0$ for $0<\alpha<1$,
the envelope critical value is positive. Then
\begin{equation}\label{eq:sec6-uniform-size}
\boxed{\displaystyle
\limsup_{n\to\infty}\sup_{\theta\in\Theta_0^T}
P_\theta\{\mathrm{LRT}_n(x_0)>
 c_\alpha^{\mathrm{pw}}(\widehat J_{V,n},x_0)\}
\le\alpha.}
\end{equation}
\end{theorem}

The uniform approximation
\eqref{eq:sec6-uniform-approx} is an additional statistical
assumption for a general marginal model, not a consequence
of fixed-parameter LAN. Appendix~\ref{app:pointwise-calibration}
proves the conditional result. For the four response families
and informative compact fitted model classes specified in
Subsection~\ref{subsec:sec4-compact-uniform}, the necessary
likelihood and moving-parameter premises are verified by
the family-specific and uniform likelihood arguments
collected in the supplementary proofs. The following result
states the consequent unconditional-in-regularity (but
model-restricted) uniform size assertion.

\subsection{Verified uniform calibration for the four GLMM families}
\label{subsec:sec6-four-family-uniform}

The preceding conditional theorem is useful for models
beyond the four families considered here. In the Gaussian
identity-link, Poisson log-link, NB2 log-link, and compound
Poisson--Gamma (CPG) log-link models, its statistical
regularity conditions can instead be established from
explicit design and compactness restrictions.

\begin{theorem}[Four-family uniform calibration]
\label{thm:sec6-four-family-uniform}
Let the independent clusters comprise a fixed finite number
of bounded deterministic design types with convergent
frequencies. Assume positive limiting frequency for an
identifying design or for the specified full-rank grouped
designs. In the basic layouts, the Gaussian family has an
identifiable fixed-effect design, at least three observations
and rank-two random-effects design; Poisson and NB2 use two
distinct covariate levels; and CPG has three conditionally
independent same-level replicates at $x=0$ and two at $x=1$.
Equal-level replicates have the same fixed, random-effect,
and family-specific predictors. The dimension of the full
nuisance vector is fixed and finite, and the identifying
design has full column rank for that parameterization.

Let the nuisance parameter and covariance ranges satisfy
the buffered compact restrictions of
\eqref{eq:sec4-compact-classes}, with exact global marginal
likelihood maximization over the enlarged fitted domain
$\Theta^+$. The Gaussian residual variance and NB2 size
remain in compact positive intervals. For CPG, both
$\phi$ and the power parameter $\kappa$ are unknown,
with $\phi$ in a compact positive interval and
$\kappa$ in a compact subinterval of $(1,2)$.
Fix $x_0\in\mathbb R$, $0<\alpha<1$, and
$\Theta_0^T\subset\mathcal C$ as in
Theorem~\ref{thm:sec6-uniform-size}; in particular,
$T\,\|w_0\|^2<M$ so that its upper covariance boundary
is inactive in the fitted model. Let
$\widetilde\lambda_n$ be the nuisance component of the
pointwise-null global maximizer, and evaluate
$\widehat J_{V,n}=J(\widetilde\lambda_n,0)$ as
\emph{model expected} efficient information at the vertex.
Then
\begin{equation}\label{eq:sec6-four-family-size}
\boxed{\displaystyle
\limsup_{n\to\infty}\ \sup_{\theta\in\Theta_0^T}
P_\theta\!\left\{
\mathrm{LRT}_n(x_0)>
 c_{\alpha}^{\mathrm{pw}}(\widehat J_{V,n},x_0)
\right\}\leq\alpha.}
\end{equation}
The conclusion permits $\Gamma=0$, fixed nonzero
rank-one null covariances, and all intermediate null
sequences within $\Theta_0^T$. It does not assume that
the CPG dispersion or power is known.
\end{theorem}

\begin{proof}
The informative designs, admissible compact ranges,
Gaussian covariance heat-derivative extension, and
family-specific moment envelopes establish common
likelihood differentiation and information identities,
identification, and uniformly positive-definite full
information. Compact-domain global likelihood localization
and a true-parameter-centered quadratic expansion then
yield uniform root-$n$ localization of both constrained
maximizers, a uniform score central limit theorem, and
consistent nuisance and information estimation.

For an arbitrary sequence of true null parameters
$\theta_n=(\lambda_n,t_n r_0)\in\Theta_0^T$, take a
subsequence along which $\lambda_n\to\lambda_*$ and
$\sqrt n\,t_n\to\tau\in[0,\infty]$.
If $\tau<\infty$, then $t_n\to0$ and the exact moving
local covariance sets converge to $K-\tau r_0$ and
$K_0(x_0)-\tau r_0$, with efficient information
$J(\lambda_*,0)$. After Fisher whitening this is the
finite-displacement cone-versus-ray Gaussian experiment;
Theorem~\ref{thm:sec6-endpoint} bounds its rejection
tail by the two-endpoint envelope.
If $\tau=\infty$, the moving sets converge to the
rank-one supporting half-space and null line. The
likelihood-ratio limit is
$\tfrac12\chi_1^2+\tfrac12\chi_2^2$, regardless of
whether $t_n$ converges to a positive value or approaches
zero more slowly than $n^{-1/2}$. Its upper-tail critical
value is no larger than
$c_{\alpha}^{\mathrm{pw}}(J(\lambda_*,0),x_0)$.
Uniform nuisance consistency, continuity of the
positive envelope critical value, and Gaussian
anti-concentration at positive thresholds justify
replacement by $\widehat J_{V,n}$. Every subsequential
rejection probability has limit superior at most
$\alpha$, which proves \eqref{eq:sec6-four-family-size}
by contradiction.
\end{proof}

\paragraph{Information evaluation and implementation scope.}
The vertex matrix in \eqref{eq:sec6-vertex-information}
is expected information evaluated under the model with
covariance set to zero and nuisance parameter
$\widetilde\lambda_n$. A sample observed Hessian at a
\emph{false} zero covariance need not estimate that matrix
when the true null covariance is nonzero. A separately
consistent estimate of efficient information at the true
parameter, obtained from a consistent full or pointwise-null
fit, may also be used with the corresponding endpoint
calibration. The theorem concerns the exact marginal
likelihood and mathematically specified critical values.
Numerically fitted log-likelihood suprema must approximate
the exact full and restricted global suprema with uniform
$o_P(1)$ errors, and numerical critical-value errors must
also vanish uniformly, before the same size guarantee
can be claimed for a software implementation.
